\section{The architectural collapse}
\label{sec:collapse}

This section states and proves the two results on which the whole design
rests. Both are \status{T}: they rest on established theorems of UHM
together with one elementary Diophantine fact and Hurwitz's theorem. Nothing
here is postulated.

\subsection{The single incidence structure}

Let $\mathsf{N}$ denote the incidence of the Fano plane $\Fano$: seven
points $\{A,S,D,L,E,O,U\}$, seven lines
\[
\{A,S,L\},\ \{A,D,E\},\ \{A,O,U\},\ \{S,D,O\},\ \{S,E,U\},\ \{D,L,U\},\ \{L,E,O\},
\]
each line carrying three points, each point lying on three lines, and every
pair of points lying on exactly one common line (the Steiner system
$S(2,3,7)$). Write $k^\ast(i,j)$ for the \emph{mediator} of a pair --- the
unique third point of their common line.

\begin{theorembox}[Theorem~\ref{thm:collapse} (Architectural collapse)]
On a seven-mode UHM substrate the four architectural layers --- interconnect,
error-correcting code, instruction algebra, and learning rule --- are the same
object: each is a function of the single incidence $\mathsf{N}$.
\end{theorembox}

\begin{theorem}\label{thm:collapse}
Let a holon's state be $\Gam\in\Dens(\mathbb{C}^7)$ evolving under the UHM
Lindbladian $\LindOmega$. Then:
\begin{enumerate}[leftmargin=1.7em,itemsep=1pt]
  \item[\textbf{(L1)}] \emph{Interconnect.} The physical coupling is the seven
        lines of $\mathsf{N}$: sector $i$ reaches sector $j$ only through the
        mediator $k^\ast(i,j)$ (the closure law: a coherence
        $\gamma_{ij}$ is created or repaired only by the live arms
        $\gamma_{i k^\ast},\gamma_{k^\ast j}$; two coherences birth a third
        iff the triple is a line).
  \item[\textbf{(L2)}] \emph{Code.} The parity-check matrix of the Hamming
        code $[7,4,3]$ has as its columns the seven nonzero vectors of
        $\mathbb{F}_2^3$; under any labelling these are the seven points of
        $\mathsf{N}$, and the sixteen codewords include exactly the seven
        weight-3 words, whose supports are the seven lines. The interconnect
        incidence \emph{is} the code's parity-check geometry.
  \item[\textbf{(L3)}] \emph{Instruction algebra.} The cubic (non-linear)
        part of $\LindOmega$ is generated by the octonionic structure
        constants $f_{ijk}$ --- the unique $G_2$-invariant trilinear form on
        $\mathrm{Im}(\Oct)$ --- which are non-zero exactly when $\{i,j,k\}$ is
        a line (the Fano selection rule). The multiplication table of the
        instruction is the line incidence of $\mathsf{N}$.
  \item[\textbf{(L4)}] \emph{Learning.} The regeneration term $\Regen$ drives
        $\Gam$ to $\rhostar=\vp(\Gam)$, and the entire flow is covariant under
        $G_2=\mathrm{Aut}(\Oct)$, whose action on the seven modes has
        $\mathsf{N}$ as its orbit incidence ($G_2$-rigidity). The fixed-point
        (learning) rule is a function of the same $\mathsf{N}$.
\end{enumerate}
Hence $\mathrm{L1}=\mathrm{L2}=\mathrm{L3}=\mathrm{L4}$ as functions of
$\mathsf{N}$. \hfill$\square$
\end{theorem}

\begin{corollary}[Zero mismatch tax]\label{cor:zerotax}
Because the four layers are one structure, there are no inter-layer
interfaces to negotiate. The area, static power, and verification effort a
conventional machine spends reconciling network~$\leftrightarrow$~code,
code~$\leftrightarrow$~ISA, and ISA~$\leftrightarrow$~optimiser
(Table~\ref{tab:four-layers}) is, in \TALOS{}, identically zero: there is
one structure to lay out, verify, and clock.
\end{corollary}

\begin{figure}[H]
\centering
\begin{tikzpicture}[font=\small,>=Stealth]
  % four separate boxes on the left
  \node[draw,rounded corners,fill=talossteel!8,minimum width=2.7cm,minimum height=0.7cm] (net)  at (0,3)   {Interconnect (NoC)};
  \node[draw,rounded corners,fill=talossteel!8,minimum width=2.7cm,minimum height=0.7cm] (ecc)  at (0,2)   {Error code (ECC)};
  \node[draw,rounded corners,fill=talossteel!8,minimum width=2.7cm,minimum height=0.7cm] (isa)  at (0,1)   {Instruction (ISA)};
  \node[draw,rounded corners,fill=talossteel!8,minimum width=2.7cm,minimum height=0.7cm] (lrn)  at (0,0)   {Learning rule};
  \node[align=center,text=plosgray] at (0,3.9) {\textbf{Today: four subsystems}\\\scriptsize designed apart, taxed at every seam};
  % the canonical Fano plane on the right (triangle + 3 medians + incircle)
  \begin{scope}[shift={(6.4,1.5)}]
    \coordinate (fA) at (0,1.7);      % A: top vertex
    \coordinate (fS) at (-1.5,-0.9);  % S: bottom-left vertex
    \coordinate (fD) at (1.5,-0.9);   % D: bottom-right vertex
    \coordinate (fL) at (-0.75,0.4);  % L: midpoint A-S
    \coordinate (fE) at (0.75,0.4);   % E: midpoint A-D
    \coordinate (fO) at (0,-0.9);     % O: midpoint S-D
    \coordinate (fU) at (0,0);        % U: centre
    % 3 edges + 3 medians (all Fano lines: {A,S,L}{A,D,E}{S,D,O}{A,O,U}{S,E,U}{D,L,U})
    \draw[talosamber,thick] (fA)--(fS) (fA)--(fD) (fS)--(fD) (fA)--(fO) (fS)--(fE) (fD)--(fL);
    % incircle through L,E,O  (line {L,E,O})
    \draw[talosamber,thick] (0,-0.034) circle[radius=0.866];
    \foreach \n in {fA,fS,fD,fL,fE,fO,fU}{ \fill[talossteel] (\n) circle (2.4pt); }
    \node[font=\scriptsize\bfseries,above] at (fA) {A};
    \node[font=\scriptsize\bfseries,left]  at (fS) {S};
    \node[font=\scriptsize\bfseries,right] at (fD) {D};
    \node[font=\scriptsize\bfseries,fill=white,inner sep=0.5pt] at (-0.95,0.5) {L};
    \node[font=\scriptsize\bfseries,fill=white,inner sep=0.5pt] at (0.95,0.5)  {E};
    \node[font=\scriptsize\bfseries,below] at (fO) {O};
    \node[font=\scriptsize\bfseries,fill=white,inner sep=0.5pt] at (0.22,0.18) {U};
  \end{scope}
  \node[align=center,text=talosamber] at (6.2,3.9) {\textbf{\TALOS{}: one object}\\\scriptsize the Fano plane $\mathsf{N}$};
  % converging arrows
  \foreach \b in {net,ecc,isa,lrn}{ \draw[->,plosgray,thick] (\b.east) to[out=0,in=180] (4.4,1.5); }
  \node[text=plosgray,font=\scriptsize] at (3.4,0.4) {collapse};
\end{tikzpicture}
\caption{Architectural monism. In every prior machine the four layers are
distinct subsystems joined at taxed interfaces (left). Theorem~\ref{thm:collapse}
proves that on a seven-mode substrate they are one incidence structure ---
the Fano plane (right) --- so the interfaces, and their tax, vanish
(Corollary~\ref{cor:zerotax}).}
\label{fig:collapse}
\end{figure}

\begin{zerobox}[From zero --- one drawing that is four games]
Chalk a hopscotch grid on the pavement. The same drawing is four different
things at once, depending on what you are doing with it. It is the
\emph{field} --- where your feet may land (the network). It is the
\emph{rule} --- which hops are legal (the instruction). It is the
\emph{referee} --- a wrong hop shows instantly, because your foot is outside
a line (the checking). And it is the \emph{teacher} --- playing the grid over
and over is what makes you better at it (the learning). Nobody bolted four
gadgets together; one chalk drawing does all four jobs \emph{because of its
shape}. Theorem~\ref{thm:collapse} says: for a seven-part self-watching
machine there exists exactly one such chalk drawing --- the seven-line figure
on the right of Figure~\ref{fig:collapse} --- and Theorem~\ref{thm:unique}
says no other number of parts has one.
\end{zerobox}

\subsection{Why exactly seven --- and nowhere else}

A skeptic will grant that \emph{if} $N=7$ then the four layers coincide, and
ask whether seven is special or merely one convenient point. It is unique,
and the reason is sharper than the consciousness-theoretic minimality proofs
of UHM: it is an arithmetic tie-break decided by the third-order principle.

\begin{theorembox}[Theorem~\ref{thm:unique} (Uniqueness of the collapse)]
The architectural collapse of Theorem~\ref{thm:collapse} is realisable at a
unique node count: $N=7$.
\end{theorembox}

\begin{theorem}\label{thm:unique}
Require three integer conditions to hold at the same $N$:
\begin{itemize}[leftmargin=1.6em,itemsep=0pt]
  \item[\textbf{(C1)}] $N$ is the order of a genuine projective plane:
        $N=q^2+q+1$ for a prime power $q\ge 2$ \emph{(so the interconnect is a
        symmetric $S(2,q{+}1,N)$ design of $N$ line-buses)};
  \item[\textbf{(C2)}] $N$ is a perfect Hamming length: $N=2^r-1$
        \emph{(so the incidence is a perfect single-error-correcting code)};
  \item[\textbf{(C3)}] $N$ is the imaginary dimension of a normed division
        algebra with a \emph{non-zero associator} \emph{(so the instruction
        algebra carries genuine third-order, non-associative structure)}.
\end{itemize}
Then $N=7$ uniquely.
\end{theorem}

\begin{proof}
\textbf{(C3) alone forces $7$.} By Hurwitz's theorem the normed division
algebras are $\mathbb{R},\mathbb{C},\mathbb{H},\Oct$ of dimensions $1,2,4,8$,
with imaginary parts of dimension $0,1,3,7$. Of these only $\Oct$ is
non-associative (its associator $[x,y,z]\ne 0$); the others are associative,
so their trilinear form is degenerate and no third-order structure exists.
Hence (C3) $\Rightarrow N=7$.

\textbf{(C1)$\cap$(C2) recurs, but only $7$ survives (C3).} The Diophantine
equation $q^2+q+1=2^r-1$ has integer solutions $q\in\{2,5,90\}$, i.e.\
$N\in\{7,31,8191\}$ (exhaustively for $q\le 200$). Of these, only $q=2,5$ are
prime powers, so only $N=7,31$ are orders of \emph{genuine} projective planes
($q=90$ is not a prime power, so $\mathrm{PG}(2,90)$ is not a Galois plane and
$N=8191$ is a bare integer coincidence, not an interconnect). Thus the
network-\emph{is}-code coincidence recurs at $N=7$ and $N=31$ --- seven is not
the only place a fabric can be a Hamming code. But
$\{7,31,8191\}\cap\{0,1,3,7\}=\{7\}$: of every coincidence, only seven is a
division-algebra imaginary dimension, and (by the paragraph above) only seven
is \emph{non-associative}. The intersection of all three conditions is
$\{7\}$. \end{proof}

\begin{remark}[The tie-break is the third-order principle]
The content of Theorem~\ref{thm:unique} is that the associator is the
discriminator. Among the node counts where the fabric can \emph{be} the code
($7,31,8191$), only seven also lets the instruction genuinely live on
\emph{triples} rather than pairs --- the UHM third-order principle, which
proves that pairwise structure is Fano-blind and that structure begins at
three. At $N=31$ or $8191$ one may still fuse interconnect and ECC, but the
instruction algebra collapses to an associative (pairwise) one and the
learning rule loses its cubic generator: the four-fold collapse fails. The
degenerate candidate $N=3$ (the quaternions) is excluded twice over ---
quaternions are associative, and $q=1$ is not a projective plane. Seven is
the only integer at which a network, a code, a non-associative instruction,
and a covariant fixed point are one thing.
\end{remark}

\begin{remark}[Relation to UHM's minimality proofs]
UHM already proves $N=7$ from consciousness (the seven-dimensionality
minimality theorem) and from algebra (the octonionic derivation and
$G_2$-rigidity). Theorem~\ref{thm:unique} is a \emph{third, independent}
route to the same integer, native to computer architecture: it asks only
that a machine's network be its code and its instruction be third-order, and
derives seven. That three unrelated demands --- a viable mind, a division
algebra, and a self-correcting third-order computer --- land on the same
number is itself the strongest evidence that seven is structural rather than
chosen. \status{T}
\end{remark}
